\chapter{Metodologia}
\label{cap:metodologia}

\section{Natureza e Abordagem da Pesquisa}
\label{sec:natureza-abordagem}

Esta pesquisa adota abordagem multim\'etodos (\textit{mixed methods}),
combinando dimens\~oes qualitativa e quantitativa em etapas sequenciais e
complementares, conforme proposto por \textcite{creswell2018} para
pesquisas aplicadas em Administra\c{c}\~ao P\'ublica. A natureza da
pesquisa \'e aplicada e de car\'ater explorat\'orio-descritivo.

A estrat\'egia de pesquisa segue a metodologia Design Science Research
(DSR), proposta por \textcite{hevner2004}, especialmente adequada para
pesquisas de Sistemas de Informa\c{c}\~ao e Administra\c{c}\~ao P\'ublica
que envolvem desenvolvimento de artefatos tecnol\'ogicos com impacto
organizacional, por exigir que o artefato resolva um problema real,
identificado empiricamente, e seja avaliado por crit\'erios de utilidade
e rigor cient\'ifico.

\section{Fundamenta\c{c}\~ao Emp\'irica: Dados e T\'ecnicas de Coleta}
\label{sec:dados-coleta}

A pesquisa utilizar\'a exclusivamente dados secund\'arios, de acesso p\'ublico, combinados a partir de fontes distintas.

\begin{itemize}
  \item A Pesquisa de Letramento Financeiro (Banco Central do Brasil em
  parceria com o Fundo Garantidor de Cr\'editos, metodologia OCDE/INFE)
  \autocite{bcbfgc2023}, que fornece evid\^encia estat\'istica
  nacionalmente representativa sobre o conhecimento da popula\c{c}\~ao
  acerca de produtos de renda fixa; o Raio X do Investidor Brasileiro
  (ANBIMA), como fonte complementar de padr\~oes de comportamento do
  investidor de varejo; e registros p\'ublicos do Reclame Aqui, filtrados
  pelo conjunto de bancos p\'ublicos comerciais de varejo do Brasil --
  Banco do Brasil e Caixa Econ\^omica Federal (federais) e BRB, Banrisul,
  Banestes, Banese e Banpar\'a (estaduais/distrital, sobreviventes do
  Programa de Incentivo \`a Redu\c{c}\~ao do Setor P\'ublico Estadual na
  Atividade Banc\'aria -- PROES) -- e por vincula\c{c}\~ao a decis\~ao
  ativa de investimento, como evid\^encia ilustrativa da manifesta\c{c}\~ao
  pr\'atica dessas barreiras. A condi\c{c}\~ao de controle p\'ublico de
  cada institui\c{c}\~ao foi verificada em agosto de 2026. Os padr\~oes de
  d\'uvidas identificados nos registros do Reclame Aqui s\~ao analisados
  via an\'alise de conte\'udo categorial \autocite{bardin2011},
  complementando a an\'alise estat\'istica descritiva da Pesquisa de
  Letramento Financeiro, e fortalecendo a validade dos dados sem exigir
  acesso direto aos clientes.

  \item S\'eries hist\'oricas extra\'idas via API do SGS/Banco Central
  (taxas de juros, \'indices de infla\c{c}\~ao, expectativas Focus);
  pre\c{c}os e taxas di\'arios do Tesouro Transparente; taxas de mercado
  da ANBIMA; normativos e resolu\c{c}\~oes da CVM e do CMN, incluindo o
  PL n\textsuperscript{o} 2.338/2023 (Marco Legal da IA) como refer\^encia
  de governan\c{c}a prospectiva. Todos os dados secund\'arios s\~ao de
  acesso p\'ublico e gratuito.
\end{itemize}

Para a etapa de avalia\c{c}\~ao (objetivo espec\'ifico e), o sistema ser\'a
testado com personas sint\'eticas -- perfis de investidor fict\'icios,
n\~ao vinculados a dados reais de clientes de nenhum banco p\'ublico -- constru\'idas em duas
camadas. A estrutura segue as tr\^es dimens\~oes de suitability da
Resolu\c{c}\~ao CVM n\textsuperscript{o} 30/2021 (objetivos de
investimento, situa\c{c}\~ao financeira e conhecimento sobre os riscos),
garantindo cobertura sistem\'atica do espa\c{c}o de teste regulat\'orio.
Os valores de calibra\c{c}\~ao de cada persona (faixa et\'aria, renda,
n\'ivel de conhecimento financeiro) s\~ao ancorados nos quatro perfis
comportamentais publicados na 8\textsuperscript{a} edi\c{c}\~ao do Raio X
do Investidor Brasileiro \autocite{anbima2025} -- Diversifica, Caderneta,
Economiza e N\~ao Investe, e Sem Reservas --, preservando a
representatividade do p\'ublico de perfil popular que \'e o foco de bancos p\'ublicos.
As fontes utilizadas no mapeamento de dores do objetivo espec\'ifico (a) --
Pesquisa de Letramento Financeiro, Raio X do Investidor e Reclame Aqui --
n\~ao s\~ao utilizadas na constru\c{c}\~ao das personas, por consistirem
em dados agregados ou de terceiros, sem perfil individual atribu\'ivel.
Essa abordagem dispensa o uso de dados reais de clientes em qualquer
etapa da pesquisa.

Cabe reconhecer explicitamente que esse desenho difere do padr\~ao de
rigor predominante na literatura de constru\c{c}\~ao de personas, que
recomenda coleta de dados prim\'arios com indiv\'iduos reais --
entrevistas em profundidade com clientes, grupos focais com
colaboradores, survey quantitativo e valida\c{c}\~ao com gestores
experientes -- como mecanismo central de rigor metodol\'ogico
\autocite{nogamisenra2025}. A op\c{c}\~ao por personas ancoradas
exclusivamente em dados p\'ublicos secund\'arios decorre da natureza
deste trabalho -- valida\c{c}\~ao de uma arquitetura t\'ecnica no
\^ambito da Administra\c{c}\~ao P\'ublica, n\~ao pesquisa de
comportamento do consumidor -- e da delimita\c{c}\~ao \'etica de n\~ao
envolver dados reais de clientes em nenhuma etapa. O mecanismo de rigor
adotado em substitui\c{c}\~ao \`a coleta prim\'aria \'e a
triangula\c{c}\~ao entre fontes p\'ublicas independentes -- detalhada a
seguir --, que cumpre fun\c{c}\~ao an\'aloga \`a triangula\c{c}\~ao de
m\'etodos qualitativos e quantitativos recomendada por
\textcite{nogamisenra2025}, com a diferen\c{c}a de operar inteiramente
sobre bases de dados secund\'arias e p\'ublicas.

\textbf{Calibra\c{c}\~ao num\'erica das personas.} A calibra\c{c}\~ao
segue um procedimento de duas etapas, aplicado uniformemente \`as
quatro personas. Primeiro, a renda de refer\^encia \'e ancorada no
ponto de maior concentra\c{c}\~ao comportamental do perfil
correspondente dentro da pr\'opria pesquisa ANBIMA -- por exemplo, para
o perfil Sem Reservas, a faixa de renda em que a concentra\c{c}\~ao
desse perfil atinge o pico --, e n\~ao em uma m\'edia agregada de
classe. Esse crit\'erio \'e adotado sempre que uma fonte de renda
independente -- as faixas do Indicador Ipea de Infla\c{c}\~ao por
Faixa de Renda \autocite{ipea2026rendafaixa}, ou a renda m\'edia por
classe da Pesquisa de Or\c{c}amentos Familiares
\autocite{ibge2019pof} -- converge para a mesma ordem de grandeza do
ponto ancorado, reforçando sua robustez por triangula\c{c}\~ao entre
fontes independentes em vez de dep\^ender de uma \'unica medi\c{c}\~ao.

Segundo, a capacidade mensal de investimento de cada persona \'e
derivada do \textit{super\'avit financeiro} da classe de rendimento
correspondente -- a diferen\c{c}a entre a renda m\'edia e a despesa
m\'edia mensal familiar da classe, ambas medidas pela Pesquisa de
Or\c{c}amentos Familiares 2017-2018 \autocite{ibge2019pof}. Essa
abordagem foi adotada em substitui\c{c}\~ao a dois indicadores
publicados diretamente pela mesma pesquisa -- \textit{varia\c{c}\~ao
patrimonial} e \textit{aumento do ativo} --, ambos descartados ap\'os
verifica\c{c}\~ao expl\'icita: o primeiro por ambiguidade de dire\c{c}\~ao
(o exemplo prim\'ario do IBGE para o conceito \'e saque de poupan\c{c}a,
sugerindo desacumula\c{c}\~ao de reservas, n\~ao investimento novo); o
segundo por cobertura incompleta, publicado apenas para as duas
classes extremas de rendimento, insuficiente para calibrar as quatro
personas. O super\'avit financeiro, calculado diretamente a partir de
duas tabelas que cobrem as sete classes de rendimento da pesquisa,
resolve essa limita\c{c}\~ao de cobertura sem depender de um indicador
pr\'e-calculado, ainda que sujeito \`a imprecis\~ao de comparar m\'edias
agregadas de m\'odulos de coleta distintos (renda e despesa) da mesma
pesquisa.

Personas cuja classe de rendimento apresenta super\'avit negativo --
caso do perfil Sem Reservas, cuja despesa m\'edia mensal supera a
renda m\'edia em aproximadamente R\$~248 -- s\~ao tratadas como sem
capacidade de investimento, e n\~ao como detentoras de um valor
residual simb\'olico. Essa n\~ao \'e uma limita\c{c}\~ao atribu\'ida
arbitrariamente \`a persona, mas o reflexo direto da condi\c{c}\~ao
financeira real dessa classe de rendimento no Brasil, conforme medida
pela pr\'opria Pesquisa de Or\c{c}amentos Familiares -- leitura
consistente com a caracteriza\c{c}\~ao do perfil na pesquisa ANBIMA
(incapacidade declarada de poupar ou investir) e com os indicadores de
estresse financeiro e endividamento j\'a registrados no diagn\'ostico
do Cap\'itulo 5.

O processo completo de verifica\c{c}\~ao de fontes -- incluindo o
levantamento das faixas de renda, a compara\c{c}\~ao entre indicadores
da POF e as corre\c{c}\~oes de precis\~ao de cita\c{c}\~ao feitas
durante a revis\~ao -- est\'a documentado de forma \'integra, com data
e justificativa de cada decis\~ao, nos arquivos de apoio do
reposit\'orio da pesquisa, permitindo auditoria externa completa da
calibra\c{c}\~ao.

\section{T\'ecnicas de An\'alise dos Dados e Matriz de Op\c{c}\~oes Metodol\'ogicas}
\label{sec:analise-matriz}

As t\'ecnicas de an\'alise variam conforme a etapa da pesquisa. Na Etapa
1, ser\~ao combinadas an\'alise estat\'istica descritiva sobre a Pesquisa
de Letramento Financeiro e an\'alise de conte\'udo categorial
\autocite{bardin2011} sobre registros do Reclame Aqui. Na fase
quantitativa de backtesting (Etapa 4), ser\~ao calculadas m\'etricas de
desempenho: retorno l\'iquido acumulado e frequ\^encia de supera\c{c}\~ao
do CDI em janelas mensais. O c\'alculo do Sharpe ratio ser\'a realizado
com a devida ressalva metodol\'ogica de que a janela de 12 meses, embora
necess\'aria para viabilizar o c\'alculo do indicador dentro do per\'iodo
dispon\'ivel, imp\~oe limita\c{c}\~oes ao poder de generaliza\c{c}\~ao
estat\'istica dos indicadores de volatilidade -- limita\c{c}\~ao que ser\'a
explicitamente reconhecida no cap\'itulo de limita\c{c}\~oes do estudo. Na
fase de avalia\c{c}\~ao (Etapa 5), ser\~ao aplicadas estat\'istica
descritiva sobre os resultados da compara\c{c}\~ao automatizada de
fidelidade e ader\^encia a suitability, e an\'alise de conte\'udo sobre
as justificativas geradas pelo LLM-as-a-judge na dimens\~ao de clareza.

\textbf{Rubrica da avalia\c{c}\~ao t\'ecnica.} Para operacionalizar H1
com maior granularidade, a avalia\c{c}\~ao t\'ecnica adotar\'a rubrica
estruturada em tr\^es pilares quantific\'aveis: (i) fidelidade do Retorno
Final L\'iquido (RFL) reportado na resposta ao cliente frente ao valor
calculado pelo m\'odulo determin\'istico -- n\~ao uma reavalia\c{c}\~ao
da matem\'atica do m\'odulo, j\'a verificada no objetivo espec\'ifico c
--, incluindo a correta aplica\c{c}\~ao da al\'iquota regressiva do IR
($\theta$) e da al\'iquota do IOF, aferida por compara\c{c}\~ao num\'erica
automatizada (desvio percentual entre o valor reportado e o valor
calculado, com limiar de toler\^ancia pr\'e-definido), sem interven\c{c}\~ao
humana; e (ii) ader\^encia da recomenda\c{c}\~ao \`as regras de
suitability vigentes (Resolu\c{c}\~ao CVM n\textsuperscript{o} 30/2021),
aferida por compara\c{c}\~ao autom\'atica contra uma matriz de adequa\c{c}\~ao
produto$\times$perfil, constru\'ida combinando as tr\^es dimens\~oes de
suitability exigidas pelo art. 3\textordmasculine{} da Resolu\c{c}\~ao CVM
n\textsuperscript{o} 30/2021 (objetivos de investimento, situa\c{c}\~ao
financeira, conhecimento de risco) com a caracteriza\c{c}\~ao de risco de
cada produto do escopo (Tesouro Selic, Prefixado, IPCA+ e CDB) -- seguindo
como precedente metodol\'ogico o desenho de benchmark de literatura em vez
de julgamento subjetivo de especialista adotado por \textcite{oehler2024},
detalhado a seguir, sem extrair diretamente faixas num\'ericas de fontes
externas que medem aloca\c{c}\~ao a\c{c}\~oes$\times$renda-fixa, dimens\~ao
fora do escopo deste trabalho. A op\c{c}\~ao por benchmark em vez de
avalia\c{c}\~ao por especialista ao vivo \'e reforçada por evid\^encia
emp\'irica de que assessores financeiros humanos frequentemente n\~ao
customizam recomenda\c{c}\~oes ao perfil declarado do cliente, mesmo
cobrando taxas elevadas \autocite{foerster2017}. A dimens\~ao (iii), clareza da
linguagem da resposta ao cliente, \'e avaliada separadamente por
LLM-as-a-judge, conforme detalhado a seguir. A rubrica permite
diagnosticar separadamente falhas de fidelidade, de conformidade
regulat\'oria e de clareza, sem exigir coleta de dados junto a
participantes humanos em nenhum dos tr\^es pilares.

\textbf{Avalia\c{c}\~ao de clareza por LLM-as-a-judge.} A valida\c{c}\~ao
da dimens\~ao de clareza da linguagem \'e realizada por meio de
LLM-as-a-judge -- uso de um modelo de linguagem para avaliar a qualidade
comunicativa das respostas geradas --, abordagem com precedente
metodol\'ogico consolidado na literatura de Processamento de Linguagem
Natural: \textcite{zheng2023}, no estudo fundador do paradigma, reportam
que ju\'izes como o GPT-4 atingem mais de 80\% de concord\^ancia com
prefer\^encia humana em tarefas abertas, patamar equivalente \`a
concord\^ancia humano-humano, resultado corroborado por \textcite{liu2023geval}
com a m\'etrica G-Eval. O julgamento \'e complementado por m\'etrica de
legibilidade validada para o portugu\^es e ancorado nas caracter\'isticas
do p\'ublico de perfil popular identificadas no diagn\'ostico do
Cap\'itulo 5 (baixa familiaridade com terminologia de renda fixa), de
modo que o LLM-juiz avalia a resposta especificamente sob o crit\'erio de
compreensibilidade para esse perfil de investidor, e n\~ao sob um padr\~ao
gen\'erico de qualidade textual.

A ado\c{c}\~ao de LLM-as-a-judge exige, contudo, reconhecimento
expl\'icito de suas limita\c{c}\~oes documentadas na literatura.
\textcite{li2024llmjudges} sistematizam os principais vieses do
paradigma: vi\'es de posi\c{c}\~ao, vi\'es de verbosidade (prefer\^encia
por respostas mais longas independentemente da qualidade) e vi\'es de
autofavorecimento (tend\^encia do modelo a avaliar mais favoravelmente
sa\'idas semelhantes \`as que ele mesmo geraria). \textcite{bavaresco2025},
em estudo de grande escala com onze modelos avaliadores em vinte tarefas
de PLN, recomendam valida\c{c}\~ao cuidadosa contra julgamento humano
antes do uso de LLM-as-a-judge como avaliador em produ\c{c}\~ao.
Reconhece-se essa recomenda\c{c}\~ao como limita\c{c}\~ao do presente
estudo: a valida\c{c}\~ao da dimens\~ao de clareza \'e realizada
integralmente por LLM-as-a-judge e m\'etrica de legibilidade, sem
valida\c{c}\~ao humana complementar dentro do escopo desta disserta\c{c}\~ao;
a implanta\c{c}\~ao em produ\c{c}\~ao da arquitetura deve, portanto, ser
precedida de uma etapa adicional de valida\c{c}\~ao humana em escala,
n\~ao coberta pelo escopo deste trabalho.

\textbf{Condi\c{c}\~ao de controle (LLM sem m\'odulo determin\'istico).}
Para verificar n\~ao apenas a acur\'acia do sistema completo, mas a
necessidade da segrega\c{c}\~ao entre c\'alculo e gera\c{c}\~ao de
linguagem proposta em H1, o mesmo conjunto de perguntas utilizado na
aferi\c{c}\~ao de fidelidade do RFL (item i da rubrica acima) ser\'a tamb\'em
submetido ao modelo de linguagem sem acesso ao m\'odulo Python de
c\'alculo, solicitando que o pr\'oprio LLM realize o c\'alculo de
rentabilidade l\'iquida -- incluindo a aplica\c{c}\~ao das al\'iquotas
regressivas de IR e IOF -- a partir dos mesmos dados de entrada. A taxa
de erro dessa condi\c{c}\~ao de controle ser\'a comparada \`a taxa de
erro da arquitetura completa, operacionalizando empiricamente a
alega\c{c}\~ao de H1 de que a segrega\c{c}\~ao entre c\'alculo
determin\'istico e gera\c{c}\~ao de linguagem elimina a principal fonte
de erro dos modelos generativos puros. O desenho segue a l\'ogica de
avalia\c{c}\~ao por remo\c{c}\~ao de componentes (\textit{ablation})
empregada em arquiteturas agentic recentes sobre RAG financeiro, nas
quais a adi\c{c}\~ao de um agente com ferramentas de c\'alculo sobre o
retrieval puro produziu o maior ganho isolado de acur\'acia factual
observado entre todas as etapas do pipeline \autocite{akinfaderin2025}.

% TODO: quando a base RAG regulatoria (objetivo especifico d) for construida,
% puxar para ca o corpus normativo ja levantado e verificado em
% bibliografia/normas-rag-corpus.md -- CVM 30/2021 (pendente de PDF), e as
% Resolucoes CMN 4.557/2017, 4.968/2021, 4.879/2020, 4.893/2021 e 5.274/2025
% (essas 5 com PDF oficial ja arquivado em bibliografia/normas/ e conteudo
% conferido diretamente). A 5.274/2025 traz a lista exata dos 14 controles
% minimos obrigatorios de seguranca cibernetica (Art. 3, Section2 da
% 4.893/2021 com redacao dada pela 5.274/2025), incluindo item XIV sobre
% monitoramento de Deep Web/Dark Web -- vale citar com precisao no texto.

\textbf{Valida\c{c}\~ao da base RAG.} A cobertura e o recall da base de
conhecimento regulat\'oria (objetivo espec\'ifico d) ser\~ao avaliados por
meio de um conjunto de teste de 30 a 50 perguntas sobre normativos da
CVM, do CMN e da ANBIMA, com respostas gabaritadas manualmente, permitindo
calcular recall@k e precis\~ao da recupera\c{c}\~ao documental antes da
integra\c{c}\~ao da base ao agente. Adicionalmente, o mecanismo de
metadados de vig\^encia descrito no objetivo espec\'ifico d ser\'a testado
com um subconjunto de perguntas envolvendo dispositivos normativos
revogados ou alterados, verificando se o sistema prioriza corretamente a
norma vigente.

A estrat\'egia de recupera\c{c}\~ao adotar\'a busca h\'ibrida (fus\~ao
ponderada entre recupera\c{c}\~ao densa e esparsa) seguida de reranking
em duas etapas, e o corpus normativo ser\'a reestruturado em formato
markdown antes da indexa\c{c}\~ao. A busca h\'ibrida e a
reestrutura\c{c}\~ao em markdown seguem \textcite{kim2025}, que
documenta ganhos de NDCG@10 decorrentes dessas duas escolhas de design
em corpora financeiros. Essas decis\~oes s\~ao refor\c{c}adas por
evid\^encia recorrente em estudos mais recentes de RAG financeiro:
\textcite{akarsu2026} mostra que a recupera\c{c}\~ao lexical (BM25)
supera a recupera\c{c}\~ao puramente densa quando o dom\'inio exige
correspond\^encia precisa de terminologia, e que a combina\c{c}\~ao
h\'ibrida seguida de reranking supera todos os m\'etodos de est\'agio
\'unico testados. O desenho espec\'ifico de reranking em duas etapas
(um reranker leve de primeira passagem seguido por um reranker de maior
capacidade sobre os candidatos filtrados) segue \textcite{lee2026}, que
reporta ganho de 6,5\% em NDCG@20 sobre uma linha de base com reranker
\'unico; \textcite{dadopoulos2025} chega a uma conclus\~ao convergente
em corpus financeiro distinto, embora com reranking de uma \'unica etapa
e outra m\'etrica (ganho de F1-score, n\~ao NDCG), refor\c{c}ando a
utilidade geral do reranking sem validar especificamente o desenho em
duas etapas. A busca h\'ibrida \'e
particularmente relevante para o corpus regulat\'orio deste trabalho, no
qual normas textualmente pr\'oximas mas normativamente distintas -- por
exemplo, a Resolu\c{c}\~ao CMN n\textsuperscript{o} 4.557 e a
Resolu\c{c}\~ao CMN n\textsuperscript{o} 4.879 -- exigem correspond\^encia
exata de termos (n\'umero da norma) que a recupera\c{c}\~ao puramente
densa, baseada em similaridade sem\^antica, tende a confundir.

Por decis\~ao de desenho, a arquitetura de recupera\c{c}\~ao n\~ao
empregar\'a expans\~ao de consulta via documento hipot\'etico (HyDE) --
t\'ecnica que, em dois estudos independentes de RAG financeiro, piorou
a qualidade de recupera\c{c}\~ao em rela\c{c}\~ao \`a busca densa
direta, atribu\'ido \`a introdu\c{c}\~ao de ru\'ido num\'erico/factual
no documento hipot\'etico gerado pelo LLM \autocite{akarsu2026,lee2026}.
Essa evid\^encia n\~ao \'e un\^anime na literatura arquivada nesta
disserta\c{c}\~ao -- \textcite{kobeissi2026}, em benchmark distinto
(FinanceBench), reporta ganhos expressivos de recall com HyDE e
Multi-HyDE --, mas optou-se pela leitura mais conservadora dado o risco
espec\'ifico de alucina\c{c}\~ao num\'erica em conte\'udo regulat\'orio,
onde precis\~ao factual \'e requisito, n\~ao apenas m\'etrica de
desempenho.

A unidade de indexa\c{c}\~ao seguir\'a a hierarquia estrutural do
pr\'oprio normativo (artigo, inciso, par\'agrafo), e n\~ao apenas
segmenta\c{c}\~ao por tamanho fixo de bloco -- abordagem an\'aloga \`a
indexa\c{c}\~ao hier\'arquica por se\c{c}\~ao/p\'agina adotada em
sistemas de RAG financeiro recentes para reduzir confus\~ao entre
trechos estruturalmente semelhantes, como normas distintas que tratam
de temas pr\'oximos \autocite{li2025fingear,choe2025hirec,kobeissi2026}.

% TODO: quando o benchmark de suitability (objetivo especifico e) for
% construido, ver bibliografia/literatura-suitability-benchmark.md.
% ATENCAO -- Jacobs/Foerster/Choi (Oehler & Horn) medem alocacao
% acoes x renda-fixa, dimensao que NAO existe no escopo desta dissertacao
% (so renda fixa: Tesouro Selic/Prefixado/IPCA+ e CDB). Servem so como
% precedente metodologico, NAO como fonte do conteudo numerico. Nao existe
% matriz oficial produto x perfil (Tesouro Nacional/ANBIMA) -- precisa ser
% CONSTRUIDA combinando CVM 30/2021 art. 3 com caracterizacao qualitativa
% de risco por produto (Selic=conservador/liquido; Prefixado=marcacao a
% mercado; IPCA+=protecao inflacao/duration alta; CDB=risco de credito).
% Para CDB acima do limite do FGC (R$250mil/CPF/conglomerado): usar
% indice de Basileia (IF.data, consulta dinamica) + rating de agencias
% quando disponivel (usar o PIOR sinal disponivel). IMPORTANTE: retirada
% de rating apos rebaixamento (caso real do BRB, dez/2025) e sinal de
% risco elevado, nao ausencia de sinal -- o sistema nao deve ter vies de
% protecao a nenhuma instituicao, inclusive o BRB. Ver Secoes 4-6 do
% arquivo para detalhe completo.
% DECISAO DE CITACAO (Secao 7 do arquivo): citar Foerster et al. (2017)
% como apoio direto ao abandono do especialista ao vivo; Jacobs et al.
% (2014) so citavel com a ressalva da Secao 4; Choi (2022) NAO citar
% (tangencial demais, nao sustenta nenhuma escolha especifica de desenho).

\textbf{Comparabilidade com a literatura.} Para viabilizar compara\c{c}\~ao
direta com estudos equivalentes, os resultados da condi\c{c}\~ao de
controle (acur\'acia da arquitetura completa vs. LLM isolado sem m\'odulo
determin\'istico) ser\~ao reportados no mesmo formato de acur\'acia
factual utilizado por \textcite{akinfaderin2025} (52,4\%
$\rightarrow$ 94,7\% no VERAFI), e as m\'etricas de cobertura e recall da
base RAG ser\~ao reportadas em NDCG@10, formato utilizado por
\textcite{kim2025} nos sete benchmarks financeiros do FinanceRAG
Challenge. A dimens\~ao de suitability \'e avaliada por benchmark
pr\'oprio de adequa\c{c}\~ao produto$\times$perfil (detalhado no par\'agrafo
anterior), seguindo como precedente metodol\'ogico o modelo de
\textcite{oehler2024} -- comparar recomenda\c{c}\~oes de IA generativa
contra um padr\~ao de refer\^encia estruturado, em vez de julgamento de
especialista ao vivo --, sem depender do conte\'udo num\'erico espec\'ifico
daquele estudo, que mede aloca\c{c}\~ao a\c{c}\~oes$\times$renda-fixa,
dimens\~ao fora do escopo desta pesquisa. Reportar os resultados nesses
formatos permite situar o desempenho desta pesquisa entre o de sistemas
equivalentes j\'a publicados, e n\~ao apenas demonstrar viabilidade
isolada da arquitetura.

% TODO: quando o modulo Python do RFL (objetivo especifico c) for escrito,
% puxar para ca a base legal das tabelas tributarias (IOF regressivo e IR
% regressivo) documentada em bibliografia/base-legal-rfl.md -- inclui as
% duas tabelas prontas para constante no codigo, a citacao normativa
% completa (Decreto no 6.306/2007 art. 32 + Anexo; Lei no 11.033/2004
% art. 1o) e a nota sobre a MP no 1.303/2025 (tentou substituir a tabela
% do IR, caducou em 11/10/2025 sem votacao).

\begin{quadro}[htbp]
  \centering
  \caption{Matriz de Op\c{c}\~oes Metodol\'ogicas}
  \label{qua:matriz-metodologica}
  \adjustbox{max width=\textwidth}{%
  \small
  \begin{tabular}{p{3.2cm}llllll}
    \toprule
    \textbf{Objetivos espec\'ificos} & \textbf{Abordagem} & \textbf{Natureza} & \textbf{Dados} & \textbf{P\'ublico-alvo} & \textbf{Coleta} & \textbf{An\'alise} \\
    \midrule
    a) Mapear barreiras informacionais & Qualitativa & Explorat\'oria & Secund\'arios (p\'ublicos online) & Survey BCB/FGC; ANBIMA; Reclame Aqui & Extra\c{c}\~ao de dados p\'ublicos & Estat\'istica descritiva e de conte\'udo \\
    b) Projetar arquitetura de agente \'unico & Qualitativa & Descritiva & Secund\'arios & Documentos regulat\'orios & An\'alise documental & S\'intese qualitativa \\
    c) Desenvolver m\'odulos Python & Quantitativa & Aplicada & Secund\'arios & APIs BCB/ANBIMA & Extra\c{c}\~ao via API & Teste de precis\~ao \\
    d) Construir base RAG & Qualitativa & Descritiva & Secund\'arios & Normativos CVM/CMN & Indexa\c{c}\~ao h\'ibrida (densa+esparsa) & Cobertura e recall \\
    e) Backtesting e avalia\c{c}\~ao & Quantitativa & Descritiva & Secund\'arios & Simuladores oficiais + benchmark de literatura + dados hist\'oricos & Compara\c{c}\~ao automatizada + simula\c{c}\~ao & Estat\'istica descritiva \\
    \bottomrule
  \end{tabular}}
  \fonte{Elaborado pelo autor (2026).}
\end{quadro}
